\section{Limitations}

The study uses one network size ($N = 144$), two noise levels, and exact recall as its main metric, which counts one wrong pixel as a failure. The synthetic families are idealized: one prototype per set, or independent biased bits. Real character sets have richer structure, such as clusters and close pairs. The loads step by one pattern from 2 to 9 and by 3 to 10 above that, so values above 9 are interpolated across wide gaps, and $P_{50}$ values between 2 and 3 are interpolated between the only two loads tested there. The bootstrap intervals on these capacities leave out the error of that interpolation. None of the classical recalls in the synthetic grid was still changing at the 50-sweep cap. The capacity here is an empirical, exact-recall, finite-$N$ measure, and it is not capacity in the sense of Amit and colleagues or of McEliece and colleagues.

The covariance rule uses the stored set's own mean activity, which is unstable for two nearly identical patterns (recall \N{r:proto:0.9:cov:10:2} at $\rho = 0.9$ with two patterns). The modern network uses one normalization and three values of $\beta$, and its $\beta$ thresholds depend on that normalization. Storkey's results depend on presentation order: in the synthetic grid, the order is the random generation order, so the intervals include it, and on the characters, the whole-set numbers are file order. The study uses Storkey's exact local field $h_{ij}$.

Near 18 stored characters, subset intervals reflect little pattern sampling, because there are only 18 distinct subsets of size 17 and one of size 18; at size 18, they measure noise, update order, and, for Storkey, presentation order. The study does not apportion the characters' failure among density, shared structure, and load, and it does not show that regularizing the pseudoinverse rule improves recall from noisy queries, because it did not vary the regularization.
